\chapter{Behaviour Problems in Teenagers}

\section*{Definition and Overview}
Behaviour problems in teenagers refer to persistent patterns of difficult, disruptive, or defiant behaviour that interfere with the young person's life at home, school, or in the community. Some conflict and moodiness are a normal part of adolescence as the brain matures and identity develops. When behaviour becomes frequent, severe, or harmful -- such as aggression, rule-breaking, truancy, or risk-taking that endangers the young person or others -- it may represent a conduct disorder or another mental health or developmental condition that warrants assessment and support.

\section*{Epidemiology}
Difficult behaviour is very common in adolescence, and most teenagers occasionally argue, test limits, or act impulsively. Significant behaviour problems affecting daily function occur in roughly 1 in 20 to 1 in 10 teenagers, with oppositional and conduct-type problems more common in boys. Rates are higher in young people with learning difficulties, ADHD, autism, family conflict or instability, and socioeconomic disadvantage. Many behavioural difficulties coexist with anxiety, depression, or substance use, and they often emerge or worsen around the early- to mid-teen years.

\section*{Aetiology and Pathophysiology}
\begin{itemize}
    \item \textbf{Brain development:} the adolescent brain matures unevenly -- the parts controlling emotion and reward develop before the areas controlling judgment and impulse control, which continue maturing into the mid-20s
    \item \textbf{Hormonal and physical changes:} puberty brings mood swings and heightened sensitivity to social cues
    \item \textbf{Genetics and temperament:} some young people inherit a more impulsive or emotionally reactive temperament
    \item \textbf{Developmental conditions:} ADHD, autism spectrum disorder, and learning difficulties are strongly associated with behaviour problems
    \item \textbf{Family and environment:} conflict at home, harsh or inconsistent parenting, family breakdown, peer pressure, and bullying all contribute
    \item \textbf{Mental health conditions:} anxiety, depression, and trauma can present as irritability, withdrawal, or acting out
    \item \textbf{Substance use:} alcohol and drugs can cause or worsen impulsive and risky behaviour
\end{itemize}

\section*{Clinical Features}
\begin{itemize}
    \item Persistent arguing, defiance, and refusal to follow reasonable rules
    \item Frequent temper outbursts, irritability, or angry moods
    \item Aggression toward family members, peers, or property
    \item Lying, stealing, or deliberate rule-breaking
    \item Truancy, declining school performance, or suspension
    \item Risk-taking: substance use, unsafe sex, dangerous driving, or running away
    \item Withdrawal from family, loss of interest in usual activities, or excessive screen use
    \item Low mood, anxiety, poor sleep, or changes in appetite
    \item Harm to self or threats to harm others
\end{itemize}

\section*{Differential Diagnosis}
\begin{itemize}
    \item \textbf{Normal adolescent development:} occasional conflict and risk-taking are expected
    \item \textbf{Depression or anxiety disorders,} which may present as irritability and withdrawal
    \item \textbf{ADHD:} inattention, impulsivity, and hyperactivity
    \item \textbf{Autism spectrum disorder:} difficulty with social rules and change can be misinterpreted as defiance
    \item \textbf{Conduct disorder or oppositional defiant disorder}
    \item \textbf{Substance use or intoxication}
    \item \textbf{Trauma or post-traumatic stress,} including after bullying or abuse
    \item \textbf{Medical causes of behaviour change,} such as thyroid disorders, head injury, or epilepsy
\end{itemize}

\section*{When to Seek Medical Care}
Seek help from a GP or school counsellor if behaviour problems are persistent, escalating, affecting school or relationships, or if you are struggling to manage them at home.

\textbf{Call 000 (or local emergency services) immediately if the teenager:}
\begin{itemize}
    \item Talks about or plans suicide, or is at immediate risk of self-harm
    \item Is threatening or has caused serious harm to others
    \item Is engaging in life-threatening risk-taking or has run away
    \item Is intoxicated, confused, or unconscious
    \item Shows sudden, severe changes in behaviour or psychosis-like symptoms (hallucinations, paranoia)
\end{itemize}

\section*{Diagnosis and Investigations}
\begin{itemize}
    \item \textbf{Detailed history:} from the teenager and parents or carers, covering behaviour at home, school, and with peers, plus development, family, and social context
    \item \textbf{Interview with the young person:} private discussion with a clinician builds trust and reveals symptoms the family may not see
    \item \textbf{Questionnaires:} standardised behaviour, mood, and ADHD rating scales filled in by parents and teachers
    \item \textbf{School information:} reports on behaviour, attendance, and learning from teachers
    \item \textbf{Developmental and cognitive assessment:} to identify learning difficulties or neurodevelopmental conditions
    \item \textbf{Mental health assessment:} evaluation for depression, anxiety, trauma, and substance use
    \item \textbf{Physical examination and tests:} to exclude medical causes, with blood tests only if clinically indicated
\end{itemize}

\section*{Management}
\begin{itemize}
    \item \textbf{Parenting and family support:} structured parent-training programs teach consistent, positive discipline and reduce conflict at home
    \item \textbf{Individual therapy:} cognitive behaviour therapy and other approaches help the young person manage emotions, solve problems, and change unhelpful thinking
    \item \textbf{School-based support:} school counsellors, tailored learning plans, and mentoring programs address truancy and academic difficulties
    \item \textbf{Behavioural programs:} structured reward systems that reinforce positive behaviour and clear consequences for harmful behaviour
    \item \textbf{Treatment of underlying conditions:} medication for ADHD, or treatment for depression, anxiety, or substance use, may be appropriate under specialist guidance
    \item \textbf{Addressing family issues:} family therapy, mediation, or social support where conflict, separation, or housing stress plays a role
    \item \textbf{Crisis services:} for severe or dangerous behaviour, youth mental health services and crisis teams may be involved
    \item \textbf{Follow-up:} regular review to monitor progress and adjust strategies as the young person develops
\end{itemize}

\section*{Complications}
\begin{itemize}
    \item School failure, suspension, or leaving school early
    \item Conflict with the law and contact with police or the justice system
    \item Substance misuse and its health and social consequences
    \item Relationship breakdown within the family and loss of friendships
    \item Self-harm and suicide -- a serious risk in distressed teenagers
    \item Long-term mental health problems, including depression and antisocial behaviour, persisting into adulthood
    \item Impact on siblings and on parents' health and wellbeing
\end{itemize}

\section*{Prognosis and Outcomes}
Many young people with behaviour problems improve significantly, especially with early and consistent support, and much of the emotional turmoil of adolescence settles as the brain matures. Oppositional and mild difficulties often resolve without specialist care. More severe conduct problems that begin early in childhood are the most persistent and are linked to poorer adult outcomes, including unemployment and mental health problems, but intervention in the teenage years still helps. The strongest predictor of a good outcome is a supportive, consistent relationship with at least one caring adult.

\section*{Prevention and Self-Care}
\begin{itemize}
    \item Maintain warm, consistent communication; listen to the teenager's views without always criticising
    \item Set clear, reasonable rules with predictable consequences, and agree them together where possible
    \item Praise positive behaviour and effort, not just results
    \item Support good sleep, regular exercise, and healthy limits on screens and social media
    \item Stay involved in school life and keep teachers informed of concerns
    \item Model calm conflict resolution and manage your own stress
    \item Encourage interests, friendships, and activities that build confidence and belonging
    \item Seek help early from a GP, school counsellor, or youth service rather than waiting for problems to escalate
    \item If you are worried about immediate safety, do not delay -- call 000 or a crisis service such as Lifeline (13 11 14)
\end{itemize}

\section*{Sources and Further Reading}
The following sources provide reliable, up-to-date information on behaviour problems in teenagers.
\begin{itemize}
    \item NHS. \textit{Mental health and behaviour in children and young people.} https://www.nhs.uk/mental-health/
    \item Mayo Clinic. \textit{Teen health.} https://www.mayoclinic.org
    \item NICE. \textit{Antisocial behaviour and conduct disorders in children and young people.} https://www.nice.org.uk
    \item World Health Organization. \textit{Adolescent health.} https://www.who.int
    \item MSD Manual. \textit{Conduct Disorder.} https://www.msdmanuals.com
\end{itemize}
